% Independent contribution. AI-produced by GPT-6 Astra; no priority claimed.
\subsection{A concrete signed-frame equalizer and the remaining fixed lift}

On the multiplier-one group $H$ retain the actual signed
R1 frames and define the finite-valued map
\[
 \Theta:H\longrightarrow
 C_{\operatorname{SL}_2(\mathbb F_3)}(R_0)
       \times C_{\operatorname{SL}_2(\mathbb F_3)}(S)
       = C_6\times C_4,
 \qquad h\longmapsto(T_h,W_h).
 \tag{B.X11}
\]
Let $P_\Theta=\Theta(H)$, without assuming that this
image is the entire displayed product.  The character
$e:P_\Theta\to\mathbb F_2$ is
$e(T,W)=\epsilon_3(T)+\epsilon_2(W)$, so
$\Delta|_H=e\Theta$.

\begin{result}{B.1.10}{Proved}{Both parities occur in the signed-frame equalizer}
The map (B.X11) is a continuous homomorphism, and
\[
 \Delta\bigl(\{h\in H:\Theta(Dh)=\Theta(h)\}\bigr)
                       =\mathbb F_2.
 \tag{B.X12}
\]
More precisely the actual full element $g$ of B.2.8 has
\[
             \Theta(g)=\Theta(Dg)=(-I,-I),
             \qquad\Delta(g)=1.
 \tag{B.X13}
\]
These equalities retain the signed matrices, rather
than only their projective images or their parity.
\end{result}
\begin{proof}
We check the composition law before using the finite
records.  For $h,h'\in H$, the unique pure frames and
the signed covariance of R1, E.65.10--E.65.12, give
\[
 \begin{gathered}
 A_{hh'}=\Phi_h(A_{h'})A_h,\quad
 \beta_{hh'}=\Phi_h(\beta_{h'})\beta_h,\quad
 V_{hh'}=V_hV_{h'},\\
 q_{12}\Phi_h=\operatorname{Ad}(V_h)q_{12}
                  \quad\hbox{on the actual pure subgroup}.
 \end{gathered}                                      \tag{B.X14}
\]
Consequently
\[
 T_{hh'}=q_{12}(A_h)^{-1}V_h
                 q_{12}(A_{h'})^{-1}V_{h'}=T_hT_{h'},
 \qquad W_{hh'}=W_hW_{h'}.
 \tag{B.X15}
\]
Continuity follows from that of the uniquely normalized
pure frames.  This proves that (B.X11) has its claimed
domain, target, and composition law.

The full canonical-cover lift of B.2.8 gives an actual
normalized $g\in H$.  Its checked equations (B.N8) and
(B.N11) are
\[
 q(\beta_g)=q(A_g)=q(\beta_{Dg})=q(A_{Dg})=[5I],
       \qquad q(V_g)=q(V_{Dg})=I.
 \tag{B.X16}
\]
The prescribed pure signed representatives are $I$
modulo four and $-I$ modulo three.  Therefore both
$T$ and $W$ equal $-I$ modulo three for both elements.
In the cubic centralizer $\epsilon_3(-I)=1$; in the
harmonic centralizer $-I=S^2$, so
$\epsilon_2(-I)=0$.  This proves (B.X13).
The identity element gives parity zero in the same
equalizer.  These are all the possible values in
$\mathbb F_2$, proving the exact image (B.X12).
\end{proof}

Here is the exact comparison with the characteristic
outer-action quotients of B.1.9.  Continuity and the
faithful compact outer action provide a characteristic
finite quotient $Q$ for which
$\ker(\rho_Q|_H)\subseteq\ker\Theta$.
There is then the uniquely specified surjective
homomorphism
\[
 \pi_Q:\rho_Q(H)\twoheadrightarrow P_\Theta,
                  \qquad \pi_Q(\rho_Q(h))=\Theta(h).
 \tag{B.X17}
\]
The kernel containment proves well-definedness.
The joint map has the exact coordinate rule
\[
 (\pi_Q,\pi_Q):
 (\rho_Q(h),\rho_Q(Dh))\longmapsto
                    (\Theta(h),\Theta(Dh)).
\]
It sends every diagonal pair to a diagonal pair.
Thus the actual
$Q$-equalizer is contained in the equalizer of
(B.X12).  No converse inclusion is used: (B.X13)
does not compute $S_m^H$ for every such $Q_m$.

\paragraph{An exact residual equation.}
Let $\Xi:H\twoheadrightarrow\prod_{n\ge0}\mathbb F_2$
be the actual character map of B.2.25 and put
$N=\ker\Xi$.  Its restriction $D|_N:N\to N$ is
surjective.  Indeed, for $n\in N$, choose $a\in H$
with $Da=n$.  Then $\Xi(a)=(e,0,0,\ldots)$.
The element $k$ of B.1.3 has $Dk=1$ and
$\Xi(k)=(1,0,0,\ldots)$; hence $ak^e\in N$
and $D(ak^e)=n$.

Choose an actual $g_0\in H$ with
$\Xi(g_0)=(1,1,\ldots)$, supplied by the proved
surjectivity, and put $a_0=g_0^{-1}Dg_0\in N$.
There is the exact equality of compact sets
\[
 \operatorname{Fix}(D)\cap\Xi^{-1}(1,1,\ldots)
 =g_0\{n\in N:a_0=n(Dn)^{-1}\}.
 \tag{B.X18}
\]
In fact every element of the fibre is uniquely $g_0n$,
and $D(g_0n)=g_0n$ is exactly $a_0Dn=n$.
Replacing $g_0$ by $g_0n_0$ changes $a_0$ to
$n_0^{-1}a_0D(n_0)$, retaining the same full lifting
problem.  The continuous map $n\mapsto n(Dn)^{-1}$
is not asserted to be a homomorphism or to be onto.
Thus (B.X18) records the remaining full equation after
all iterated quadratic coordinates have been fixed.

For the physical fixed subgroup there is an exact
full-multiplier reduction with a different arithmetic
correction.  The actual map $D_L=\operatorname{Ad}_cD$
fixes $G_{\mathbb Q}$ pointwise, and therefore
\[
 C_L=(C_L\cap H)G_{\mathbb Q},\qquad
       \Delta(C_L)=\Delta(C_L\cap H).
 \tag{B.X19}
\]
Indeed write $g=h\gamma$ with $h\in H$ and
$\gamma\in G_{\mathbb Q}$.  The equation
$D_L(g)=g$ reduces literally to $D_L(h)=h$.
Both factors on the right then belong to $C_L$;
arithmetic vanishing of $\Delta$ gives the image equality.
Moreover $D(c)=c$ implies that $D$ commutes with
$\operatorname{Ad}_c$, whence
\[
 \ker D_L=\ker D,\qquad D_L(H)=H,
 \qquad \Delta D_L^n=\Delta D^n\quad(n\ge0).
 \tag{B.X20}
\]
Thus the same actual character fibre and residual
equation (B.X18), with $D_L$ substituted for $D$,
also describe the physical multiplier-one fixed lift.
